% !PRINTED 59
\weberpage{59}

Ersetzt man $\varphi(u)$ durch $\varphi(u)-c$, worin $c$ einen beliebigen
Wert bedeutet, so folgt der wichtige Satz:

3. \gesperrt{Die Summe der Argumentwerte, in denen eine
doppeltperiodische Funktion im Inneren eines Periodenparallelogramms
 einen und denselben Wert $c$ annimmt,
ist von $c$ unabhängig oder ändert sich bei stetiger Veränderung
 von $c$ höchstens sprungweise um eine Periode.}

4. \gesperrt{Hieraus folgt, daß es auch keine doppeltperiodische
 Funktion erster Ordnung gibt.}

Denn eine solche Funktion müßte jeden Wert $c$ in einem
Punkte $u$ annehmen. Ist sie also gleich $c'$ in einem Punkte $u'$
und gleich $c''$ in einem davon verschiedenen Punkte $u''$, so würde
aus 3. $u'\equiv u''$ folgen, gegen die Voraussetzung.

Unter einer \gesperrt{doppeltperiodischen Funktion der zweiten
Art} versteht man nach \gesperrt{Hermite} eine Funktion $\psi(u)$, die den
Perioden $\omega_1$, $\omega_2$ gegenüber sich den Gleichungen gemäß verhält:

\[ \psi(u+\omega_1) = a_1\,\psi(u)\,, \quad \psi(u+\omega_2) = a_2\,\psi(u)\,, \]

worin $a_1$, $a_2$ Konstanten sind.

5. \gesperrt{Auch für eine doppeltperiodische Funktion der
zweiten Art gilt der Satz, daß sie in gleich vielen
Punkten des Periodenparallelogramms Null und unendlich
 wird.} Ist die Zahl dieser Punkte $m$, so heißt auch hier die
Funktion von der $m$ten Ordnung.

Der Beweis ist derselbe wie der des Satzes 2. Es gilt aber
hier nicht mehr die Erweiterung, daß die Funktion $\psi(u)$ jeden
Wert gleich oft annimmt, weil $\psi(u)-c$ jetzt nicht mehr periodisch
 ist.

6. \gesperrt{Eine doppeltperiodische Funktion $\psi(u)$ von der
zweiten Art und der Ordnung Null ist notwendig eine
Konstante oder eine Exponentialfunktion.} Denn wenn $\psi(u)$
nicht Null wird, so ist

\[ \varphi(u) = \dfrac{\psi'(u)}{\psi(u)} \]

eine doppeltperiodische Funktion von der ersten Art, die nicht
unendlich wird, also nach 1. eine Konstante $a$. Daraus folgt:

\[ \psi(u) = A\,e^{a\,u}\,, \]

worin $A$ eine neue Konstante ist.
